% Quantum Gravity Dynamics — new Chapter 4, independent research manuscript
% This is NOT an edit to the historical Chapter 4.
% Original 19 October 2026; QGD-097 updated 26 October 2026.
% QGD-098 surgical referee-verification insertion 26 October 2026.
% QGD-101 surgical integration of research notes QGD-099/100, 26 October 2026.
% QGD-115 update, 26 October 2026: Einstein-Hilbert form derived by uniqueness (Lovelock) under principle L.
% QGD-113 final update, 26 October 2026: curved spin-1/2 flow law and slow-motion limits.
% QGD-112 update, 26 October 2026: exact spin-1/2 flow law; naming of derived objects corrected.
% QGD-111 update, 26 October 2026: covariant spin-1/2 version of the quantum terms.
% QGD-110 insertion, 26 October 2026: quantum structure of the master metric and
% covariance tests (from research notes QGD-107--109 and the QGD-110 tests).
% QGD-105 audit correction, 26 October 2026: QGD-101 insertions restricted to what
% follows from this chapter's stated assumptions; reframing of QGD's identity removed.
% All QGD-101 insertions are marked; no prior equation or label is removed.
% Units c=hbar=1 in variational arguments.
\documentclass[11pt,a4paper]{article}
\usepackage[T1]{fontenc}
\usepackage[utf8]{inputenc}
\usepackage[margin=25mm,headheight=15pt]{geometry}
\usepackage{amsmath,amssymb,amsthm,mathtools,bm,booktabs,microtype,xcolor,mathrsfs}
\usepackage[colorlinks=true,linkcolor=blue!55!black,urlcolor=blue!55!black]{hyperref}
\usepackage{fancyhdr}
\pagestyle{fancy}\fancyhf{}\lhead{Quantum Gravity Dynamics}\rhead{Chapter 4 — Natural-Invariant Dynamics}\cfoot{\thepage}
\newtheorem{theorem}{Theorem}[section]
\newtheorem{proposition}[theorem]{Proposition}
\newtheorem{remark}[theorem]{Remark}
\newcommand{\dd}{\mathrm d}
\newcommand{\cE}{\mathcal E}
\newcommand{\cR}{\mathcal R}
\newcommand{\cG}{\mathscr G}
\newcommand{\cT}{\mathscr T}
\newcommand{\cQ}{\mathscr Q}
\newcommand{\cS}{\mathscr S}
\title{\textbf{Quantum Gravity Dynamics}\\[6pt]
\large Chapter 4: Natural-Invariant Dynamics\\ and the Physical Meaning of Gravity\\[5pt]
\normalsize QGD--098 research edition --- referee-verified radial Dirac audit and semiclassical benchmark}
\author{QGD research manuscript --- independent of the archived historical Chapter 4}
\date{26 October 2026 --- QGD--101 integration, QGD--105 audit correction, QGD--110 quantum structure, QGD--111/112/113 spin-$\frac12$ updates, QGD--115 Lovelock}
\begin{document}
\maketitle
\begin{abstract}
What does the natural-invariant description of gravitation add to the
Einsteinian metric equations? We formulate a precise first answer, without
postulating a new wave law. Given the composite QGD metric of Chapters 1--2
and the effective gravitational and matter action used in Chapters 2--3,
variation with respect to the independent natural-invariant one-forms
produces an exact projected Euler equation. The historical three-channel
idea thereby admits an unambiguous \emph{action-defined successor},
$\cS=\cG+\cT+\cQ=0$. In the derivative-free, metric-only sector,
$\cQ=0$ and the remaining terms are projections of the Einstein--matter
residual. On a regular full-rank signed-field representation, those
projected equations are equivalent to the effective Einstein equations.
The result clarifies the physical role of $\sigma_\mu$, the origin of its
representation freedom, and the exact location where genuinely distinct
QGD dynamics would have to arise. We incorporate the Chapter~3 breakthroughs:
Newtonian calibration of the exterior natural-invariant amplitude; the
exact static Killing-energy expansion; a corrected conserved-current
bridge; the leading Dirac-derived quantum density-gradient stress; and a
no-go theorem for universal power-law exterior corrections produced solely
by restructuring a compact spherical source at fixed mass. It does not assert a microscopic
wave equation, a derived $\sigma$ kinetic term, or a new gravitational
observable. This short chapter is intended to grow as these questions
are solved; it does not alter the historical Chapter 4. The chapter proves a metric-only quantum-equivalence theorem, imports a
verified semiclassical massless-Dirac far-field correction, and audits the
historical $9/2$ force. A QGD--110 section derives the quantum structure of the
master metric from the founding principles: a universal quantum sector fixed by the
source's quantum stress, a per-wave quantum conformal factor, and the test of the
historical derivative term, and the exact flow law of spin-$\frac12$ matter in a curved
universal metric; it also records Lorentz-covariance and gauge tests of the main
equations. A QGD--115 subsection shows that the Einstein--Hilbert form of the
gravitational dynamics is not an independent axiom: it follows uniquely, by
Lovelock's theorem, from QGD's metric-only structure and a second-order locality
principle, with the coupling fixed by the Newtonian calibration. The revision history is in
Section~\ref{sec:revision-history}.
\end{abstract}

\noindent\textbf{Reading key.} \textbf{D}: derived from stated assumptions;
\textbf{A}: adopted effective dynamics; \textbf{C}: fixed definition or
normalization; \textbf{P}: proposed physical interpretation not yet an
independent empirical result. These labels distinguish logical status,
not scientific importance.

\section{From the quantum description to effective geometry}
\label{sec:motivation}
General relativity establishes how geometry behaves when coupled to
energy and momentum. It does not, by itself, establish whether geometry
is fundamental or an effective macroscopic description. QGD investigates
a concrete alternative interpretation: the gravitational field is
associated with macroscopic quantum structure of matter, represented by
gravitational natural-invariant covectors $\sigma_\mu^{(s)}$; their
composites define an effective metric (P). The word \emph{natural
invariant} is QGD's historical name for these covectors; a covector is
not itself a coordinate-invariant scalar.

The essential distinction is between an \emph{ontological} proposal and a
\emph{dynamical} result. The first three chapters supply (i) the
four-component macroscopic Dirac structure and controlled local spinor
ratios, (ii) the master metric and its signed-sector representations,
and (iii) Einsteinian effective dynamics and Newtonian calibration in
specified limits. None of these statements alone proves that the
individual natural-invariant components constitute independent
observable particles or waves.

The general QGD metric is retained unchanged:
\begin{equation}\label{eq:master}
\boxed{g_{\mu\nu}=T^a{}_{\mu}T^b{}_{\nu}
\left[\eta_{ab}-\sum_{s=1}^{N}\varepsilon_s
\sigma_a^{(s)}\sigma_b^{(s)}-\lambda\mathcal Q_{ab}\right],
\quad \varepsilon_s=\pm1.}
\end{equation}
Here $T^a{}_{\mu}$ is the reference coframe, and the derivative
operator $\mathcal Q_{ab}$ and coefficient $\lambda$ remain
unmatched proposals in the unrestricted theory. We will derive definite
equations in the explicitly stated fixed-coframe, $\lambda=0$,
metric-only effective sector and then explain what changes when the
full metric map is used.

\section{The new Chapter 4 field equation}
\label{sec:balance}
In this section set $c=1$. Restore the physical Einstein coupling
$8\pi G_N/c^4$ when using SI units. Consider the composite effective
action (A)
\begin{equation}\label{eq:action}
\Gamma[\sigma,\Psi]=S_{\rm EH}[g(\sigma)]
+S_{\rm matter}[g(\sigma),\Psi],\qquad
S_{\rm EH}=\frac{1}{16\pi G_N}
\int\dd^4x\sqrt{-g}\,(R-2\Lambda).
\end{equation}
Boundary variations are compactly supported or supplied with the
appropriate gravitational boundary term. The matter equations are
varied independently. Define $T^{\mu\nu}$ by
$\delta S_{\rm matter}=\frac12\int\sqrt{-g}\,
T^{\mu\nu}\delta g_{\mu\nu}$ for on-shell matter variations, and let
\begin{equation}\label{eq:Eg}
\boxed{\cE_g^{\mu\nu}:=\frac1{\sqrt{-g}}
\frac{\delta\Gamma}{\delta g_{\mu\nu}}
=-\frac{1}{16\pi G_N}
\left(G^{\mu\nu}+\Lambda g^{\mu\nu}
-8\pi G_N T^{\mu\nu}\right).}
\end{equation}
This is a definition of the effective metric Euler derivative,
not an independent QGD field equation.

% ===================== QGD-115 INSERTION =====================
\subsection{Why the Einstein--Hilbert form: a uniqueness theorem (QGD--115)}
\label{sec:lovelock}
The action~\eqref{eq:action} was introduced as an adopted effective action (A).
Its \emph{form} need not be adopted. It follows from facts established within QGD
together with one structural principle.

\paragraph{Established within QGD.}
\begin{enumerate}
\item[(i)] \emph{Matter couples only through $g$.} The microscopic action uses a
vierbein of $g$ and is locally Lorentz invariant, so matter depends on the natural
invariants only through $g$ (Proposition~\ref{prop:nobeyond}).
\item[(ii)] \emph{Covariance.} The master metric is tensorial under coordinate
changes and invariant under frame Lorentz transformations
(Section~\ref{sec:quantumstructure}).
\item[(iii)] \emph{Conservation.} Diffeomorphism invariance of the matter action
gives $\nabla_\mu T^{\mu\nu}=0$ on shell.
\item[(iv)] \emph{No consistent dynamics beyond $g$ has been found.} The routes
tested, physical frame redundancies, an antisymmetric pairing field and two
dynamical frames, were either inconsistent at two-derivative order or required an
independent postulate (Section~\ref{sec:physics}).
\end{enumerate}

\paragraph{Principle L (locality and order).} The gravitational equation for the
natural invariants is local, generally covariant, depends on $\sigma$ only through
$g$, and contains at most second derivatives of $g$.

\begin{theorem}[Einstein form by uniqueness (D$\mid$L; Lovelock)]\label{thm:lovelock}
Under (i)--(iii) and Principle~L, in four spacetime dimensions the gravitational
equation has the form $\mathcal E_{\mu\nu}[g]=\kappa\,T_{\mu\nu}$ with
$\mathcal E_{\mu\nu}$ symmetric, built from $g$ and its first two derivatives, and
identically divergence free. By Lovelock's theorem \cite{Lovelock1971,Lovelock1972}
the only such tensors are $aG_{\mu\nu}+bg_{\mu\nu}$, so
\begin{equation}\label{eq:lovelock}
\boxed{G_{\mu\nu}+\Lambda g_{\mu\nu}=\kappa\,T_{\mu\nu}.}
\end{equation}
The Newtonian calibration of Chapter~3 then fixes $\kappa=8\pi G_N/c^4$, and the
action~\eqref{eq:action} is the corresponding variational form, up to boundary
terms.
\end{theorem}
\begin{proof}
The equation must be divergence free because its source is, by~(iii), for every
matter configuration. Lovelock's theorem classifies the symmetric, divergence-free
concomitants of $g$ with at most second derivatives in four dimensions. The
coupling is the single constant left, and F5 fixes it.
\end{proof}
\noindent\textbf{Verification.} On a generic metric with all ten components
varying, a 50-digit computation gives $\nabla^\mu G_{\mu\nu}=O(10^{-14})$, while
$\nabla^\mu R_{\mu\nu}$ and $\nabla^\mu(Rg_{\mu\nu})$ are of order $0.1$--$0.7$; the
unique divergence-free combination $aR_{\mu\nu}+bRg_{\mu\nu}$ has $b/a=-0.5000$.

\paragraph{Status.} The \emph{form} of the gravitational dynamics is derived
(D$\mid$L), and the \emph{coupling} is fixed by the calibration F5. Three things are
not derived:
\begin{itemize}
\item \textbf{$\Lambda$.} The theorem allows a cosmological term but does not fix its
value.
\item \textbf{The second-order requirement.} It is a principle, chosen for a
well-posed, ghost-free classical limit. Higher-derivative terms, such as the
logarithmic curvature-squared terms induced by quantum matter, are treated as
subleading corrections.
\item \textbf{The value of $G_N$ from the microscopic theory.} Two routes remain
open. In the induced-gravity route (QGD--114), the one-loop effective action of
Chapter~1's fields yields the Einstein form, but with one spin-$\frac12$ field and
one Maxwell field the induced Newton constant is negative and the induced
cosmological constant is of cutoff order \cite{Sakharov,Visser}. In the
thermodynamic route \cite{Jacobson}, horizon entropy counts physical degrees of
freedom positively and $\Lambda$ appears as an integration constant, but an
equilibrium hypothesis and a physical cutoff are required.
\end{itemize}
% =================== END QGD-115 INSERTION ===================

Write the metric-only signed specialization locally as
\begin{equation}\label{eq:signed}
g_{\mu\nu}=\bar g_{\mu\nu}
+\sum_s s_s\sigma^{(s)}_\mu\sigma^{(s)}_\nu,
\qquad s_s:=-\varepsilon_s\in\{+1,-1\}.
\end{equation}
All variations below treat the covectors as independent components
in a fixed frame and hold the background coframe fixed.

\begin{theorem}[Action-defined natural-invariant equation (D$\mid$A)]
\label{thm:projection}
For the action \eqref{eq:action} and metric map \eqref{eq:signed},
the complete Euler derivative of each independent natural-invariant
covector satisfies
\begin{equation}\label{eq:Euler}
\boxed{\cS_{(s)}^\rho
:=\frac1{\sqrt{-g}}\frac{\delta\Gamma}
{\delta\sigma^{(s)}_\rho}
=2s_s\cE_g^{\rho\nu}\sigma^{(s)}_\nu=0.}
\end{equation}
Equivalently, with \emph{action-defined} channels,
\begin{equation}\label{eq:split}
\boxed{\cS_{(s)}^\rho
=\underbrace{-\frac{s_s}{8\pi G_N}
\left(G^{\rho\nu}+\Lambda g^{\rho\nu}\right)
\sigma^{(s)}_\nu}_{\displaystyle\cG_{(s)}^\rho}
+\underbrace{s_s T^{\rho\nu}\sigma^{(s)}_\nu}
_{\displaystyle\cT_{(s)}^\rho}
+\underbrace{0}_{\displaystyle\cQ_{(s)}^\rho}=0.}
\end{equation}
\end{theorem}
\begin{proof}
Vary the signed metric while holding other sectors fixed:
$\delta g_{\mu\nu}=s_s(\delta\sigma^{(s)}_\mu
\sigma^{(s)}_\nu+\sigma^{(s)}_\mu\delta\sigma^{(s)}_\nu)$.
Symmetry of $\cE_g^{\mu\nu}$ gives
$\delta\Gamma=\int\sqrt{-g}\,2s_s\cE_g^{\rho\nu}
\sigma^{(s)}_\nu\delta\sigma^{(s)}_\rho$.
Equation \eqref{eq:Euler} follows from independent arbitrary
compactly supported variations. Substituting \eqref{eq:Eg}
proves \eqref{eq:split}.
\end{proof}

\begin{remark}[What the three symbols mean]
$\cG$ is the effective \emph{geometric-action contribution}, $\cT$
is the \emph{matter-stress contribution}, and $\cQ$ denotes additional
terms arising from a \emph{specified} derivative-dependent metric map
when such a map has actually been defined and varied. Thus $\cQ=0$
in the metric-only algebraic sector. The letter $Q$ does \emph{not}
by itself mean ``quantum''. In particular, a genuine quantum
matter-loop correction to a metric-only effective action modifies
$\cE_g$, hence the geometric/matter balance, without necessarily
producing any separate $\cQ$ channel.
\end{remark}


\section{What changes in the unrestricted derivative-dependent metric map}
\label{sec:derivativemap}
The master proposal~\eqref{eq:master} permits a nonzero term
$-\lambda\mathcal Q_{ab}$. Its functional form and coefficient are
currently unmatched; the existence of a symbol $\mathcal Q_{ab}$
does not establish the independent field equation $\square_g\sigma_\mu$
or any particular force correction. A conditional chain rule can,
however, be proved \emph{before} choosing that operator.

\begin{proposition}[Functional chain rule for a first-derivative metric map]
\label{prop:jet}
Hold the reference coframe fixed and suppose the composite metric is a
smooth local function $g_{\mu\nu}(\sigma,\partial\sigma)$ of the
independent natural-invariant fields and their first coordinate
derivatives. Assume the matter and gravitational action depend on
$\sigma$ only through the resulting metric, and its metric Euler
variation is
$\delta\Gamma=\int\dd^4x\,\sqrt{-g}\,\mathcal E_g^{\mu\nu}\,
\delta g_{\mu\nu}$ after the appropriate boundary treatment.
Then, for arbitrary compactly supported $\delta\sigma$, the
natural-invariant Euler derivative obeys
\begin{equation}\label{eq:jet}
\boxed{\displaystyle
\frac{\delta\Gamma}{\delta\sigma^{(s)}_\rho}
=\sqrt{-g}\,\mathcal E_g^{\mu\nu}
\frac{\partial g_{\mu\nu}}{\partial\sigma^{(s)}_\rho}
-\partial_\alpha\!\left(
\sqrt{-g}\,\mathcal E_g^{\mu\nu}
\frac{\partial g_{\mu\nu}}
{\partial(\partial_\alpha\sigma^{(s)}_\rho)}\right).}
\end{equation}
If $g$ depends on second or higher derivatives of $\sigma$, the
corresponding alternating higher total derivatives are required.
Direct nonmetric $\sigma$ couplings add their \emph{own} Euler
functional derivatives.
\end{proposition}
\begin{proof}
Use the ordinary chain rule on the specified local metric map,
substitute into the metric Euler variation, and integrate by parts
the term proportional to
$\partial_\alpha\delta\sigma^{(s)}_\rho$. Compact support removes
the boundary contribution. A second integration-by-parts argument
handles second derivatives, and so on.
\end{proof}
This equation expresses \emph{how} a well-defined derivative metric
operator would enter; it does not determine an independent kinetic
coefficient from the already accepted action. Without the actual
$\mathcal Q_{ab}$ and its admissible boundary terms, no numerical
quantum-force coefficient can be assigned to this contribution.

% === QGD-101 insertion, corrected QGD-105 ===
\begin{remark}[The derivative channel on Chapter~2's constructed branch (QGD--101/105)]
\label{rem:Q-absorbed}
For both trial completions of $\mathcal Q$ in Chapter~2 and any $\lambda\neq0$,
Chapter~2 constructs a particular nine-sector branch on which the metric map has
a differential right inverse (Theorems on the nine-sector lift). On that branch,
a metric-mediated action gives exactly the effective metric equations, whatever
$\lambda$ is. As Chapter~2 states, this concerns a particular spanning branch,
not all field configurations; on other branches, or with $\sigma$ dependence not
passing through $g$, the question remains open.
\end{remark}
% === end QGD-101/105 insertion ===

\section{Does this equation say anything beyond GR?}
\label{sec:equivalence}
One projected field equation per covector need not contain all ten
independent symmetric metric equations. What matters is the rank of
the map from independently variable signed covectors to the ten
metric components.

\begin{proposition}[Exact regular signed section (D)]
\label{prop:section}
In a fixed local reference frame let $h_{\mu\nu}
=g_{\mu\nu}-\bar g_{\mu\nu}$ be any real symmetric $4\times4$
matrix. For any constant $a>0$ set
\begin{equation}\label{eq:PN}
\boxed{P=aI_4+\frac{h}{4a},\qquad
N=aI_4-\frac{h}{4a}.}
\end{equation}
The four rows of $P$ and four rows of $N$ define eight signed
covectors, and
\begin{equation}\label{eq:identity}
\boxed{P^{\mathsf T}P-N^{\mathsf T}N=h.}
\end{equation}
For every symmetric variation $\delta h$, the choice
$\delta P=\delta h/(4a)$ and $\delta N=-\delta h/(4a)$
yields $\delta g=\delta h$. Thus the metric differential has rank ten
on this section.
\end{proposition}
\begin{proof}
Expand both products in \eqref{eq:identity}; the $a^2I$
and $h^2/(16a^2)$ terms cancel, leaving $h$.
Differentiation gives the displayed right inverse.
\end{proof}

\begin{theorem}[Conditional equivalence with Einsteinian dynamics (D$\mid$A)]
\label{thm:equiv}
Suppose the full effective action depends on the signed covectors
\emph{only through} $g(\sigma)$, and the admissible covector
variations possess a local right inverse to arbitrary symmetric
metric variations, as on the section \eqref{eq:PN}. Then
\begin{equation}\label{eq:equivalence}
\boxed{\cS_{(s)}^\rho=0\ \text{for every sector }s
\quad\Longleftrightarrow\quad
\cE_g^{\mu\nu}=0.}
\end{equation}
Thus, on this regular full-rank branch, the action-derived
natural-invariant equations are equivalent to the \emph{specified}
effective Einstein--matter equations.
\end{theorem}
\begin{proof}
If $\cE_g=0$, every projected Euler derivative is zero.
Conversely, let every covector Euler derivative vanish.
Use the differential right inverse of Proposition~\ref{prop:section}
to express any compactly supported symmetric metric test variation
as an admissible signed-covector variation. Then
$\delta\Gamma=0$ for all metric tests and the fundamental lemma
of the calculus of variations gives $\cE_g=0$.
\end{proof}

\noindent\textbf{Interpretation.}
Writing the Einsteinian equations using natural invariants is a
nontrivial local representation and an alternative choice of variables.
It does \emph{not}, by itself, generate additional observables beyond
GR. In the metric-only sector the transformation
$S\mapsto O(x)S$ with $S=\binom{P}{N}$ and
$O^{\mathsf T}\mathrm{diag}(I_4,-I_4)O=\mathrm{diag}(I_4,-I_4)$
leaves $g$ unchanged. The resulting local $O(4,4)$ freedom is a
parametrization redundancy wherever it is a symmetry of the complete
action and matter couplings. A physical $\sigma$ representative need
not be unique; adding direct interactions could change this symmetry.

\section{The calibrated gravitational field and the normalization distinction}
\label{sec:connection}
Chapter~3 independently varies both components $A(r)$ and $B(r)$ of
the declared classical spherical one-form $\sigma=A\,\dd x^0+B\,\dd r$,
without setting $A=B$ in the action. On its nondegenerate, asymptotically
flat vacuum branch, the adopted Einstein--Hilbert sector gives
\begin{equation}\label{eq:classical}
\boxed{\displaystyle A(r)^2=B(r)^2=\frac{C_g}{r},\qquad
 C_g=\frac{2G_NM}{c^2}=:r_s.}
\end{equation}
The value of $C_g$ is fixed \emph{once} by the Newtonian exterior
acceleration. It is not derived from the microscopic Dirac equation;
its calibration is legitimate correspondence with the observed
Newtonian mass parameter, as in the historical Einstein normalization.
No other coefficient is adjusted to recover higher orders.

When an independently selected reference fermion mass $m_{\rm ref}>0$
supplies a reduced Compton length $\bar\lambda_{\rm ref}=\hbar/(m_{\rm ref}c)$,
Chapter~3 defines the dimensionless \emph{geometric inverse-strength
convention}
\begin{equation}\label{eq:alphageom}
\boxed{\displaystyle |\alpha_G^{\rm geom}|^2:=\frac{\bar\lambda_{\rm ref}}{C_g}
=\frac{\hbar c}{2G_NMm_{\rm ref}}.}
\end{equation}
The original common spinor multiplier $\alpha_G^{\rm spin}$ in
$\Psi=(\alpha_G^{\rm spin}/r)\sum_a\psi_a u_a e^{\pm iS_a/\hbar}$
is \emph{not} thereby normalized from the geometric equations.
Indeed, the simultaneous change
$\alpha_G^{\rm spin}\mapsto z\alpha_G^{\rm spin}$,
$\psi_a\mapsto\psi_a/z$ with $z\ne0$ leaves the physical state
unchanged. Its absolute normalization comes from an independently
specified current or occupied-state convention. In particular,
$\alpha_G^{\rm spin}=\alpha_G^{\rm geom}$ is not a theorem of
Chapters~1--3.

The observable static Killing-energy potential per unit held test
mass is
\begin{equation}\label{eq:pot}
\boxed{\displaystyle\Phi_{\rm stat}(r)=c^2\left(\sqrt{1-\frac{C_g}{r}}-1\right),
\qquad r>C_g.}
\end{equation}
Its convergent, wholly \emph{classical} expansion is
\begin{equation}\label{eq:classseries}
\boxed{\displaystyle\Phi_{\rm stat}(r)=-\frac{G_NM}{r}
-\frac{G_N^2M^2}{2c^2r^2}
-\frac{G_N^3M^3}{2c^4r^3}
-\frac{5G_N^4M^4}{8c^6r^4}+\cdots .}
\end{equation}
All powers of $\hbar$ cancel when \eqref{eq:alphageom} is used
merely as a parametrization of $C_g$. Thus these higher orders are
not evidence of quantum-corrected vacuum gravity.

\section{The microscopic Dirac bridge: what it fixes}\label{sec:bridge}
For one locally on-shell, positive-energy, outward-propagating Dirac
WKB branch with conserved particle flux $\dot N>0$, Chapter~1 and
Chapter~3 give the magnitude-only current identity
\begin{equation}\label{eq:fluxbridge}
\boxed{\displaystyle |\alpha_G^{\rm spin}\psi_0(r)|^2
=\frac{\dot N\,\Delta(r)}{8\pi c^2p(r)},\qquad
\Delta=E_{\rm loc}+m_{\rm eff}c^2,\quad p>0.}
\end{equation}
This is a transport constraint. It determines the envelope only
relative to an independently fixed momentum profile and flux; it does
\emph{not} determine a $1/r$ gravitational potential, the geometric
normalization, or an absolute phase. A static bound source with
zero net radial current is \emph{not} normalized using this outgoing
flux formula. Its spinor profiles must be obtained from its own
coupled boundary-value equations with conserved occupation.

\section{A controlled quantum matter contribution to the source}\label{sec:quantumsource}
To make the $\hbar$ dependence precise without silently inserting a
new gravitational operator, use a stated positive-energy,
fixed-spin, nonrelativistic Dirac-to-Pauli limit and write
$\varphi=\sqrt n\,e^{iS/\hbar}\chi$ with constant unit Pauli
spinor $\chi$, $n>0$, and fermion mass $m_f>0$.
The polar identity for the Schr\"odinger kinetic energy is
\begin{equation}\label{eq:quantumenergy}
\boxed{\displaystyle E_{\rm kin}[n,S]
=\int\dd^3x\left[
\frac{n|\bm\nabla S|^2}{2m_f}
+\frac{\hbar^2}{8m_f}\frac{|\bm\nabla n|^2}{n}\right].}
\end{equation}
The quantum density-gradient functional and associated momentum-flux
tensor (one admissible improvement convention) satisfy
\begin{equation}\label{eq:qstress}
\boxed{\displaystyle Q[n]=-\frac{\hbar^2}{2m_f}
\frac{\nabla^2\sqrt n}{\sqrt n},\qquad
\Pi^Q_{ij}=\frac{\hbar^2}{4m_f}
\left(\frac{\partial_i n\partial_j n}{n}-\partial_i\partial_j n\right),\qquad
\partial_j\Pi_{ij}^Q=n\partial_i Q.}
\end{equation}
This is a derived coefficient of \emph{quantum matter stress},
not a derived coefficient of a standalone
$\hbar^2(\nabla\sigma)^2$ gravitational kinetic term.
To leading nonrelativistic order, the \emph{adopted} sourced
Einstein--Hilbert equations reduce to Poisson's law:
\begin{equation}\label{eq:Uquantum}
\boxed{\displaystyle
\nabla^2\Phi=4\pi G_Nm_fn,\qquad
\mathcal U:=-\frac{2\Phi}{c^2},\qquad
\nabla^2\mathcal U=-\frac{8\pi G_Nm_f}{c^2}n.}
\end{equation}
If the physical source profile $n(x;\hbar)$ is independently determined
from its quantum dynamics, its $\hbar$ dependence influences the
sourced metric and consequently the corresponding temporal
natural-invariant \emph{combination}. This is a source effect, not a
new vacuum propagation law. Relativistic stress, pressure, spin,
multi-orbital occupations and binding require the complete
Dirac--Einstein source instead of simply substituting one $n$.

\begin{proposition}[Compact spherical exterior at fixed mass (D$\mid$A)]
\label{prop:exterior}
Let $\rho_\hbar(r)$ be an integrable, nonnegative spherically symmetric
mass density with fixed total mass
$M=4\pi\int_0^\infty\rho_\hbar(r)r^2\,\dd r$, compactly supported
in $r\leq R$. Assume Poisson's equation and $\Phi(\infty)=0$.
Then $\Phi(r)=-G_NM/r$ and
$F_r=-G_NMm_t/r^2$ for every $r>R$, independently of the interior
$\hbar$ dependence.
\end{proposition}
\begin{proof}
Spherical Poisson integration gives
$r^2\Phi'(r)=4\pi G_N\int_0^r\rho_\hbar(s)s^2\,\dd s$.
The integral equals $M/(4\pi)$ outside the support. Integrate once
more with the stated asymptotic boundary condition.
\end{proof}
Thus a universal exterior term proportional to
$G_NM\hbar^2/(m_f^2c^2r^4)$ does \emph{not} follow from changing only
the profile of a compact spherical quantum source while keeping its
mass and the Poisson law fixed. Noncompact sources, altered total mass,
non-spherical multipoles or independently derived quantum field
equations are distinct questions.
% === QGD-113 insertion ===
\begin{remark}[Varying spin (QGD--113)]
The fixed-spin assumption above is a special case. When the spin direction
$\hat s$ varies, the exact spin-$\frac12$ flow law of Section~\ref{sec:spinhalf} adds the
spin quantum potential $(\hbar^2/8m_f)|\nabla\hat s|^2$ and a spin-vorticity force; both
vanish for constant $\hat s$, recovering Eq.~\eqref{eq:qstress}.
\end{remark}
% === end QGD-113 insertion ===

\section{Status of the historical quantum-force candidate}\label{sec:historicalforce}
The proposed universal attractive force
\begin{equation}\label{eq:historic9}
\boxed{\displaystyle F_{r,\rm hist}=-\frac{G_NMm_t}{r^2}
\left[1+\frac92\left(\frac{\hbar}{m_fc r}\right)^2+\cdots\right]}
\end{equation}
remains a historical \emph{candidate}, not a conclusion of the
accepted action. The reciprocal of $e^{2im_fcr/\hbar}$ is
$e^{-2im_fcr/\hbar}$: inverting nonconstant Taylor terms separately
does not invert the complete exponential. The historical cubic energy
also gives a sign opposite to an enhancement of attraction.

A first-order weak-field expansion of a specified Liouville-normal,
separated Schwarzschild radial Dirac equation produces coefficients
depending on its angular channel $\kappa$. Their derivatives are
not coordinate-, tetrad-, or field-redefinition-invariant force
observables. In particular the lowest $\kappa=-1$ channel does not
produce the historical term at order $G_N\hbar^2/r^4$ in that normal
form. Nothing here establishes a universal quantum change to the
exterior gravitational acceleration. To claim one, first derive a
covariant $\hbar$-dependent action, vary it including all connection
and metric dependence, and calculate a gauge-/representation-invariant
physical observable.



\subsection{Independent referee test of the radial Dirac coefficient (QGD--098)}
\label{sec:radial-referee}
The following is a \emph{conditional theorem about a quantum test field}.
It does not modify the natural-invariant field equations. Adopt the
classical, Newton-calibrated Schwarzschild exterior of
Eq.~\eqref{eq:classical}, minimally couple a massive four-component
Dirac field through the induced orthonormal-frame spin connection,
and use the standard separated, fixed-tetrad radial basis. This
additional coupling and the reduction to a single radial mode are
assumptions of the \emph{test-field model}, not consequences of the
flat-space spinor-component ratios alone. Keep $r>r_s>0$ and define
\begin{equation}\label{eq:radial-parameters}
\ell_C:=\frac{\hbar}{m_fc},\quad x:=\frac r{\ell_C},\quad
 a:=\frac{r_s}{\ell_C},\quad e:=\frac{E}{m_fc^2},\quad
 F:=1-\frac ax,\quad \frac{\dd}{\dd x_*}:=F\frac{\dd}{\dd x},
\end{equation}
where $\kappa\in\mathbb Z\setminus\{0\}$ is the spinor angular
index. For the stated radial spinor gauge the first-order field
equations read
\begin{equation}\label{eq:radial-first-order}
\boxed{\begin{aligned}
\frac{\dd f}{\dd x_*}&=(e+\sqrt F)g-\frac{\kappa\sqrt F}{x}f,\\
\frac{\dd g}{\dd x_*}&=\frac{\kappa\sqrt F}{x}g+(\sqrt F-e)f.
\end{aligned}}
\end{equation}
\begin{proposition}[Liouville-normal radial coefficient (D$\mid$A)]
\label{prop:radial-test}
Assume Eqs.~\eqref{eq:radial-parameters}--\eqref{eq:radial-first-order},
$e=1$, and $a/x<1$. Put $B:=1+\sqrt F$, $K:=\kappa\sqrt F/x$, and
\begin{equation}\label{eq:radial-pq}
P:=\frac{F_x}{F}-\frac{B_x}{B},\qquad
Q:=F K_x-F\frac{B_x}{B}K-K^2+1-F.
\end{equation}
Then elimination of $g$ gives
\begin{equation}\label{eq:radial-second-order}
f_{xx}+P f_x+\frac Q{F^2}f=0.
\end{equation}
After the invertible local Liouville rescaling
$f=\exp(-\tfrac12\int P\,\dd x)u$, the transformed equation is
\begin{equation}\label{eq:radial-liouville}
\boxed{u_{xx}+\mathcal V_\kappa(x,a)u=0,\qquad
\mathcal V_\kappa=\frac Q{F^2}-\frac12P_x-\frac14P^2.}
\end{equation}
Its expansion at fixed $x>0$ about $a=0$ is
\begin{equation}\label{eq:radial-series-general}
\boxed{\begin{aligned}
\mathcal V_\kappa={}&-\frac{\kappa(\kappa+1)}{x^2}
+\frac ax+\frac{3-4\kappa^2-\kappa}{4}\frac a{x^3}\\
&+\frac{2a^2}{x^2}
+\frac{-64\kappa^2+4\kappa+69}{64}\frac{a^2}{x^4}
+O(a^3).
\end{aligned}}
\end{equation}
In particular,
\begin{equation}\label{eq:radial-kminus1}
\boxed{\mathcal V_{-1}(x,a)=\frac ax+\frac{2a^2}{x^2}
+\frac{a^2}{64x^4}+O(a^3),}
\end{equation}
so the $O(a/x^3)$ correction vanishes in the $\kappa=-1$
channel. The displayed expansions are local analytic expansions
under the stated hypotheses; they are not uniform near $r=r_s$.
\end{proposition}
\begin{proof}
Solve the first equation in \eqref{eq:radial-first-order} for
$g=(Ff_x+Kf)/B$. Differentiate with respect to $x_*$, substitute the
second equation, and collect the $f_x$ and $f$ coefficients to
obtain Eq.~\eqref{eq:radial-second-order}. Direct differentiation
of the Liouville rescaling gives Eq.~\eqref{eq:radial-liouville}.
Using $\sqrt{1-a/x}=1-a/(2x)-a^2/(8x^2)+O(a^3)$ in the resulting
\emph{exact} coefficient gives Eq.~\eqref{eq:radial-series-general}.
Setting $\kappa=-1$ proves Eq.~\eqref{eq:radial-kminus1}.
\end{proof}
\begin{remark}[What is and is not a quantum force]
\label{rem:radial-not-force}
At $e=1$ one may \emph{formally} rewrite
Eq.~\eqref{eq:radial-liouville} as a zero-energy one-dimensional
Schr\"odinger equation, with
$U_{\rm radial}:=-m_fc^2\mathcal V_\kappa/2$.
Its $O(a/x^3)$ contribution is
\begin{equation}\label{eq:radial-gamma}
\delta U_\kappa=-\gamma_\kappa
\frac{G_NM\hbar^2}{m_fc^2r^3},\qquad
\gamma_\kappa:=\frac{3-4\kappa^2-\kappa}{4}.
\end{equation}
Differentiating this \emph{representation-dependent radial term} yields
$-3\gamma_\kappa G_NM\hbar^2/(m_fc^2r^4)$, but this is
\textbf{not} by itself a local force observable or an invariant
correction to the gravitational metric. In particular,
$\gamma_{-1}=0$ and $\gamma_{+1}=-1/2$. The historical 
\emph{attractive} coefficient $9/2$ would require
$\gamma_\kappa=3/2$, i.e.
$4\kappa^2+\kappa+3=0$, which has no real root.
The $a^2/(64x^4)$ term at $\kappa=-1$ is likewise an
\emph{effective radial-equation coefficient}, not a new vacuum
Einstein tensor. It must not be reported as a distinct QGD
quantum-gravitational force.
\end{remark}
\paragraph{Referee/proof boundary.}
The result above rests on a stated classical Schwarzschild exterior,
the minimally coupled test Dirac equation, a fixed tetrad and spinor
angular basis, and a specified threshold energy $e=1$. It therefore
follows from Chapters~1--2 only \emph{after} their declared effective
metric/Dirac coupling is adopted. It is not derived solely from the
flat-space macroscopic Dirac ratios. Exact symbolic differentiation
verifies Eqs.~\eqref{eq:radial-second-order}--\eqref{eq:radial-kminus1};
this is not a Lean-kernel proof. A physical quantum-gravity claim
would additionally require the modified, conserved metric/source
field equations and an invariant observable, with no further fit
to the historical $9/2$ coefficient.


\section{A published semiclassical Dirac benchmark in natural-invariant variables}
\label{sec:diracbenchmark}
The historical $9/2$ Compton-scale force is not established by the
original phase manipulation.  There is, however, a distinct and
quantitatively fixed $\hbar$-dependent \emph{gravitational} correction
that can be imported into the effective metric sector.  Carlson,
Anderson, Fabbri, Fagnocchi, Hirsch and Klyap computed the leading
far-field renormalized stress tensor of one free \emph{massless
four-component} Dirac field in an asymptotically empty (Boulware)
state, and checked its gravitational backreaction against previous
one-loop calculations~\cite{Carlson2010}.  This is a benchmark of
semiclassical Einstein--Dirac physics, not an independent derivation
of gravitational dynamics from the original flat-space QGD Dirac action.

\begin{proposition}[Massless Dirac far-field gravitational correction
(D$\mid$A)]\label{prop:diracfarfield}
Adopt the metric-only effective Einstein--Hilbert action with a
minimally coupled quantized free massless four-component Dirac field,
the renormalized Boulware state, asymptotic flatness, fixed ADM mass
$M>0$, areal radius $r$, and leading-order semiclassical backreaction.
In units $G_N=c=\hbar=1$, the published leading stress tensor is
\begin{equation}\label{eq:diracRSET}
\boxed{\left\langle T^\mu{}_{\nu}\right\rangle_{\rm ren}
=\frac{M}{40\pi^2r^5}\,
\operatorname{diag}(4,2,-3,-3)+O(r^{-6}).}
\end{equation}
For the far-field static metric
\begin{equation}\label{eq:scmetric}
\dd s^2=-f(r)\dd t^2+
\left(1-\frac{2\mathfrak m(r)}r\right)^{-1}\dd r^2
+r^2\dd\Omega^2,\qquad
f(r)=\left(1-\frac{2\mathfrak m(r)}r\right)e^{2\nu(r)},
\end{equation}
the asymptotically normalized functions satisfy
\begin{align}
\mathfrak m(r)&=M\left(1+\frac1{5\pi r^2}\right)
+O(r^{-3}),\label{eq:massfn}\\
\nu(r)&=\frac{M}{15\pi r^3}+O(r^{-4}),\label{eq:nufn}\\
f(r)&=1-\frac{2M}{r}-\frac{4M}{15\pi r^3}
+O(r^{-4}),\label{eq:scf}
\end{align}
where in the last expression $O(r^{-4})$ also includes mixed
classical--semiclassical orders.  Restoring $G_N,c,\hbar$, the
\emph{leading far-field} supported-static-body Killing-energy
potential per unit test mass is
\begin{equation}\label{eq:QGD97Phi}
\boxed{\begin{aligned}
\Phi_{\rm QGD}^{\rm sc}(r)={}&c^2\left(
\sqrt{1-\frac{2G_NM}{c^2r}}-1\right)\\
&-\frac{2G_N^2M\hbar}{15\pi c^3r^3}
+O\!\left(\frac{G_N^3M^2\hbar}{c^5r^4}\right)
+O(\hbar^2),
\end{aligned}}
\end{equation}
where the final $O(\hbar^2)$ denotes omitted higher quantum orders,
not a uniform SI-dimensional bound.  The displayed quantum coefficient
is meaningful in the far field $r\gg r_s=2G_NM/c^2$, where the
quantum backreaction is perturbatively small.  It must not be
extrapolated unchanged to $r\simeq r_s$.
\end{proposition}
\begin{proof}
For the ansatz \eqref{eq:scmetric}, the radial Einstein equations
at leading quantum order read
\begin{equation}\label{eq:scEFE}
\mathfrak m'=-4\pi r^2\langle T^t{}_t\rangle_{\rm ren},
\qquad
\nu'=\frac{4\pi r
(\langle T^r{}_r\rangle_{\rm ren}
-\langle T^t{}_t\rangle_{\rm ren})}
{1-2\mathfrak m/r}.
\end{equation}
Substituting \eqref{eq:diracRSET} gives
$\mathfrak m'=-2M/(5\pi r^3)$ and
$\nu'=-M/(5\pi r^4)+O(r^{-5})$.
Integrating with $\mathfrak m(\infty)=M$ and $\nu(\infty)=0$
proves \eqref{eq:massfn}--\eqref{eq:nufn}.  Consequently
$f=1-2M/r-4M/(15\pi r^3)+O(r^{-4})$.
For an observer held static in this exterior,
$\Phi=c^2(\sqrt{f}-1)$; expanding to first order in the
renormalized Dirac stress and restoring dimensions proves
\eqref{eq:QGD97Phi}.  This is a sourced asymptotic expansion,
not the solution of a freely chosen scalar force law.
\end{proof}

\begin{remark}[Quantum gravitational force and Planck-scale residue]
The \emph{quantum part} of the signed radial force on a held static
test mass $m_t$ has leading far-field behavior
\begin{equation}\label{eq:scforce}
\boxed{\delta F_r^{\rm sc}
=-\frac{2G_N^2Mm_t\hbar}{5\pi c^3r^4}
+\text{higher orders}.}
\end{equation}
Relative to the Newtonian attraction, its leading magnitude is
$\frac{2}{5\pi}\frac{G_N\hbar}{c^3r^2}$,
which is governed by the Planck length squared, not the inverse
square of an arbitrarily selected probe mass.  It is therefore
\emph{not} the historical $(9/2)(\hbar/(m_fc r))^2$ correction.
The classical piece $-m_t\Phi_{\rm stat}'$ remains defined by the
exact background square root.  The quantum term is a controlled
far-field correction, not an exact all-radius force.
\end{remark}

\begin{remark}[Natural-invariant reconstruction and status]
The full signed/tetrad QGD master metric \eqref{eq:master} is
unchanged.  For the nondegenerate one-sector static spherical
\emph{metric representation} discussed in Chapter~2, write
$\sigma=A(r)\dd t+B(r)\dd r$ and put
$F=1-A^2$, $D=1+B^2-A^2$ so that the diagonal radial metric
coefficient is $D/F$.  Reconstruct the semiclassical metric by
$A^2=1-f$ and $D/F=(1-2\mathfrak m/r)^{-1}$.
Then, in the far field and to the first displayed quantum order,
\begin{equation}\label{eq:scsigmaAB}
\boxed{\begin{aligned}
A^2(r)&=\frac{2M}{r}+\frac{4M}{15\pi r^3}
+O(r^{-4}),\\
B^2(r)&=\frac{2M}{r}+\frac{2M}{5\pi r^3}
+O(r^{-4})
\end{aligned}}
\qquad(G_N=c=\hbar=1).
\end{equation}
Thus the familiar equality $A^2=B^2$ is a property of the
classical Schwarzschild branch, \emph{not} of the quantum-corrected
source.  This construction shows that QGD variables can
represent the published semiclassical metric; it does not
establish a new dynamical law for $\sigma_\mu$.
\end{remark}

\paragraph{Assumptions and formal-verification boundary.}
The algebraic implications from \eqref{eq:diracRSET} and
\eqref{eq:scEFE} to \eqref{eq:QGD97Phi} are suitable for a
Lean-style theorem once analytic asymptotic bounds and positivity
conditions on the metric are specified.  The renormalized-stress
coefficient in \eqref{eq:diracRSET} is a \emph{published quantum
field-theoretic input}, not a theorem derived from Chapter~1.
The complete original four-mode wavefunction has not been proved
to generate this Boulware vacuum as its unique quantum state.
At the current metric-only effective level, the equivalence theorem
in the next section applies: this correction is already predicted
by semiclassical GR with the same Dirac matter.


\section{Can QGD predict a quantum force independently of semiclassical GR?}
\label{sec:independent}
This is the central discriminating question of QGD--096. Here
\emph{semiclassical Einstein gravity} means a classical metric sourced by
a renormalized expectation value or effective action of quantum matter;
\emph{quantized gravitational effective field theory} includes
fluctuations of the gravitational field itself. The two comparisons are
different. A Dirac wave in a prescribed metric, the quantum-potential
term of a Pauli reduction, and an $\hbar$-dependent source profile are
not, by themselves, quantum corrections to the gravitational vacuum
response. The relevant test is an invariant observable calculated from
an effective dynamics which QGD derives but which the selected GR-plus-
quantum-matter comparator does not already possess.

\begin{theorem}[Metric-only quantum-equivalence obstruction (D$\mid$A)]
\label{thm:quantumno-go}
Let $\mathcal F$ be the algebraic signed-sector map $\sigma\mapsto
g[\sigma]$ on a regular branch whose differential is surjective onto
symmetric metric variations, with identical boundary conditions and
matter states on both sides. Let $\Gamma_{\rm sc}[g,\Psi]$ be a
specified semiclassical Einstein--matter effective functional,
possibly including renormalized quantum matter loops, and assume
\begin{equation}\label{eq:qgd-sc-pullback}
\boxed{\Gamma_{\rm QGD}[\sigma,\Psi]
=\Gamma_{\rm sc}[\mathcal F(\sigma),\Psi]}
\end{equation}
with no independent $\sigma$ interaction or separately quantized
$\sigma$ measure. Then stationary QGD fields on this branch and
stationary fields of $\Gamma_{\rm sc}$ determine the same set of
physical metrics and metric/matter observables. In particular, no
additional universal $\hbar$-dependent gravitational-force coefficient
is implied solely by expressing $g$ through $\sigma$.
\end{theorem}
\begin{proof}
For any admissible variation, the functional chain rule gives
\begin{equation}\label{eq:quantum-pullback-chain}
\delta\Gamma_{\rm QGD}
=\int\dd^4x\,\sqrt{-g}\,\mathcal E_{\rm sc}^{\mu\nu}
\frac{\partial g_{\mu\nu}}{\partial\sigma^{(s)}_\rho}
\delta\sigma^{(s)}_\rho
+\delta_\Psi\Gamma_{\rm sc},
\end{equation}
where repeated sector indices are summed. A stationary metric of
$\Gamma_{\rm sc}$ therefore makes the QGD pullback stationary.
Conversely, by surjectivity every compactly supported symmetric
$\delta g$ is induced by an admissible $\delta\sigma$; stationarity
of the pullback and the independently imposed matter equations imply
$\mathcal E_{\rm sc}^{\mu\nu}=0$. Because all specified observables
are functionals of the same physical metric and matter state, they
agree. The result is an on-shell, mean-effective-action theorem: it
makes no claim about quantizing the redundant variables with a new
functional measure, about rank-deficient branches, or about a
new action not covered by \eqref{eq:qgd-sc-pullback}.
\end{proof}

\begin{remark}[What the theorem does, and does not, rule out]
It rules out \emph{new observable coefficients arising from a metric
reparametrization alone}. It does not forbid a distinct QGD quantum
mechanism. For example, an independently matched covariant correction
$\Delta\Gamma_{\rm QGD}[\sigma,\Psi;\hbar]$ could alter the field
equations, whether through additional natural-invariant physics, a
quantum-derived metric operator, or a consistently specified quantum
measure. But a coefficient copied from a known GR effective action,
or chosen to reproduce the historical force, would not constitute an
independent QGD prediction. A meaningful comparison must specify
the same physical source, asymptotic conditions, renormalization
prescription, and gauge-invariant observables.
\end{remark}

\subsection{An exact exterior test of the historical $9/2$ law}
\label{sec:poisson-residual}
The attractive force candidate in Eq.~\eqref{eq:historic9} can be
integrated unambiguously at the displayed order. Put
$\ell_C=\hbar/(m_fc)$, where the reference species $m_f$ must be
specified independently. In the weak-field nonrelativistic
interpretation, $F_r=-m_t\partial_r\Phi$ and $\Phi(\infty)=0$ imply
\begin{equation}\label{eq:quantum-candidate-potential}
\boxed{\displaystyle
\Phi_{9/2}(r)=-\frac{G_NM}{r}
-\frac{3G_NM\ell_C^2}{2r^3}
+\text{higher-order terms}.}
\end{equation}
For the isolated displayed $r^{-3}$ term and $r>0$,
$\nabla^2(r^{-3})=6r^{-5}$, so
\begin{equation}\label{eq:quantum-laplacian-test}
\boxed{\displaystyle
\nabla^2\Phi_{9/2}(r)=-\frac{9G_NM\ell_C^2}{r^5}\ne0
\quad (r>R),}
\end{equation}
where $R$ is the support radius of an assumed compact source and
higher-order terms have been omitted consistently. If one tried to
retain the \emph{unmodified} Poisson equation, the same radial tail
would demand an exterior effective density
\begin{equation}\label{eq:quantum-rho-eff}
\boxed{\displaystyle
\rho_{\rm eff}(r)=-\frac{9M\ell_C^2}{4\pi r^5},
\qquad r>R,}
\end{equation}
which is not the original empty exterior. This is a diagnostic
\emph{effective} density, not a claim about negative physical matter.
Therefore the historical universal correction cannot be generated
solely by changing a compact spherical source profile under the same
vacuum Poisson equation; it requires a justified exterior quantum
source/response, a modified field equation, or other physics.

At the fully relativistic level, this weak-field calculation is not
by itself a complete covariant metric solution. The independently
varied Einstein--Hilbert vacuum branch of Chapter~2 already fixes
the exact spherical exterior to Schwarzschild. A new universal
$1/r^3$ metric perturbation must consequently be checked against the
\emph{complete} modified field equations and their conservation and
boundary conditions; a scalar force ansatz cannot perform that check.
If $m_f$ denotes the test particle mass, its Compton scale also makes
the force state/species-dependent, so the claim is not automatically a
universal exterior metric correction.

\subsection{A falsifiable QGD-specific matching test}
\label{sec:delta-action}
A future new term should be expressed as the difference between a
microscopically justified QGD effective functional and the specified
semiclassical comparator:
\begin{equation}\label{eq:delta-gamma}
\boxed{\displaystyle
\Delta\Gamma[\sigma,\Psi;\hbar]
:=\Gamma_{\rm QGD}[\sigma,\Psi;\hbar]
-\Gamma_{\rm sc}[g[\sigma],\Psi;\hbar].}
\end{equation}
At fixed physical matter state and boundary data, a nonzero
$\Delta\Gamma$ is \emph{necessary but not sufficient} for a new
observable: it may contain a boundary term, a redundancy, or an
operator removable by permissible redefinitions. The minimum
verification gates are:
\begin{enumerate}
\item Derive $\Delta\Gamma$ from the stated microscopic/macroscopic
 Dirac construction or declare any additional physical axiom.
\item Specify its covariant dependence on the full signed master
 metric, $\sigma^{(s)}$, and any derivative operators; include its
 independent normalization, units and boundary terms.
\item Vary the \emph{complete} functional, including connection,
 coframe, stress-energy and possible nonlocal contributions, and
 check Bianchi/conservation identities and regular branches.
\item Fix at most the permitted Newtonian calibration coefficient,
 then predict a second observable without re-fitting it.
\item Compare with \emph{both} semiclassical GR plus quantum matter
 and, where applicable, low-energy quantized-GR effective-field-theory
 corrections, at matched sources and renormalization convention.
\end{enumerate}
A quantum Dirac radial equation in the unmodified Schwarzschild metric
meets none of the first three gates for a new \emph{gravitational
field} operator. The previously derived quantum density-gradient
stress meets the microscopic matter-side test, but its gravitational
consequences within Eq.~\eqref{eq:action} are also consequences of
the chosen semiclassical Einstein--matter system. Hence the present
QGD action does \emph{not yet} close the deeper foundational gap of
demonstrating that QGD itself predicts a quantum correction
independently of semiclassical general relativity.


% ===================== QGD-110 INSERTION =====================
\section{The quantum structure of the master metric (QGD--110)}\label{sec:quantumstructure}
This section uses only the founding principles: F1, the microscopic Dirac field;
F2, coarse-graining and the macroscopic spinor equation; F3, the master metric
with its reference term; F4, the adopted Einstein--Hilbert-plus-matter action;
F5, the Newtonian calibration. Imported results are marked L.

\subsection{A metric-level quantum term and the field equations}
\begin{proposition}[Exact split identity (D)]\label{prop:split}
Write the physical metric as $g=g_{\rm cl}+q$. If $G[g]=8\pi G\,T$, then
identically
\begin{equation}\label{eq:split}
G[g_{\rm cl}]=8\pi G\,(T+T_Q),\qquad T_Q:=-\frac{G[g]-G[g_{\rm cl}]}{8\pi G}.
\end{equation}
\end{proposition}
\noindent A quantum term in the master metric therefore always produces a
quantum-corrected Einstein equation for the classical part. By itself this is a
rewriting. It acquires content only when F1--F5 fix $q$, which they do through the
quantum stress of the matter.

\subsection{The quantum term in the static natural-invariant variables}
\begin{proposition}[Exact identity (D)]\label{prop:ABFH}
In the static branch ($F=1-A^2$, $D=1+B^2-A^2$, areal radius),
\begin{equation}\label{eq:ABFH}
A^2-B^2=1-FH,\qquad \frac{\mathrm d}{\mathrm dr}\ln(FH)=\frac{8\pi r(\rho+p_r)}{1-2\mathfrak m/r},
\end{equation}
with $H=D/F$. Hence $A=B$ exactly whenever $\rho+p_r=0$: in vacuum, for Maxwell
fields and for a cosmological constant.
\end{proposition}
\noindent For the vacuum of the microscopic Dirac field (Proposition~\ref{prop:diracfarfield}),
Eq.~\eqref{eq:scsigmaAB} gives
\begin{equation}\label{eq:ABquantum}
A^2-B^2=-\frac{2}{15\pi}\,\frac{G_N^2M\hbar}{c^5r^3},
\end{equation}
because vacuum polarization of the microscopic Dirac field has $\rho+p_r<0$. Inserting this $q$ into
Eq.~\eqref{eq:split} returns $T_Q=-T^{\rm Dirac}$ in all three independent
components (verified symbolically). This is the quantum term F1 and F4 place in the
master metric. Its scale is $\ell_P^2$, not the historical $\ell_Q$, and by
Theorem~\ref{thm:quantumno-go} it coincides with semiclassical GR.

\begin{proposition}[Test of the historical derivative trial (D)]\label{prop:trialtest}
For the classical exterior sector $A=B=\sqrt{2M/r}$ on the flat spherical
reference, the trial $\mathcal Q_{ab}=(\bar\nabla_a\sigma^c)(\bar\nabla_b\sigma_c)$ has
only angular components, $\mathcal Q_{\theta\theta}=2M/r$, and in areal gauge gives,
to first order in $M$, $\Delta f=0$ and $\Delta h=-4\lambda M/R^3$. The derived term
requires $\Delta f=-4M\ell_P^2/(15\pi R^3)$ and $\Delta h=+2M\ell_P^2/(5\pi R^3)$.
No $\lambda$ matches both, so the trial is not the quantum term of the foundations.
\end{proposition}

\subsection{A quantum conformal factor derived from the wave equation}
\begin{proposition}[Quantum conformal factor (D)]\label{prop:qcf}
For a scalar wave $\psi=Re^{iS/\hbar}$ of mass $M$ obeying
$(\square_g-M^2c^2/\hbar^2)\psi=0$, the real part of the wave equation is exactly the
Hamilton--Jacobi equation of a conformally rescaled metric,
\begin{equation}\label{eq:qcf}
\tilde g^{\mu\nu}\partial_\mu S\partial_\nu S+M^2c^2=0,\qquad
\tilde g_{\mu\nu}=\Omega_\psi^2g_{\mu\nu},\qquad
\Omega_\psi^2=1-\lambda_C^2\frac{\square_gR}{R},\quad\lambda_C=\frac\hbar{Mc}.
\end{equation}
Wherever $\Omega_\psi^2>0$, the integral curves of $\tilde g^{\mu\nu}\partial_\nu S$
are timelike geodesics of $\tilde g$.
\end{proposition}
\begin{proof}
The real part is $g^{\mu\nu}\partial S\partial S+M^2c^2=\hbar^2\square_gR/R$ (verified
symbolically). For the second statement, $u_\mu=\partial_\mu S/(Mc)$ is a gradient
with $\tilde g(u,u)=-1$, so $u^\nu\tilde\nabla_\nu u_\mu=\frac12\partial_\mu\tilde g(u,u)=0$.
\end{proof}
\noindent This is the quantum potential of Chapter~1 (Eq.~wkb-quantumHJ) and of
Section~\ref{sec:quantumsource}, in geometric form. Calling the integral curves
\emph{trajectories} is an interpretation. The theorem itself is a statement about
the wave's phase and amplitude.

\subsection{How the quantum term enters the master metric}
\begin{theorem}[Two consistent routes, one inconsistent (D$\mid$A)]\label{thm:routes}
Under F1--F5:
\begin{enumerate}
\item \textbf{Per wave (consistent).} $\Omega_\psi^2$ multiplies the universal
master metric for that wave (spin-0 matter; spin-$\frac12$ matter in
Section~\ref{sec:spinhalf}),
$\tilde g^{(\psi)}_{\mu\nu}=\Omega_\psi^2T^a{}_\mu T^b{}_\nu[\eta_{ab}+\sum_ss_s\sigma^{(s)}_a\sigma^{(s)}_b]$,
and governs that wave's phase dynamics. Universal free fall is recovered as
$\lambda_C\to0$; at finite $\lambda_C$ the mass dependence is that of ordinary
quantum mechanics in a gravitational field.
\item \textbf{Through the source's stress tensor (consistent, universal).} On
shell the wave's stress tensor splits exactly as
\begin{equation}\label{eq:TQ}
T_{\mu\nu}=\rho\,u_\mu u_\nu+T^Q_{\mu\nu},\qquad
T^Q_{\mu\nu}=\hbar^2\bigl[\partial_\mu R\partial_\nu R-\tfrac12g_{\mu\nu}\bigl((\partial R)^2+R\,\square_gR\bigr)\bigr],
\end{equation}
with $\rho=M^2c^2R^2$ and the on-shell $u_\mu=\partial_\mu S/(Mc)$. Under F4, $T^Q$
curves spacetime with coupling $8\pi G$, and through Eq.~\eqref{eq:split} it
appears as a universal quantum sector $q$ of the master metric. For a coherent
static source of size $L$, $T^Q_{00}/\rho c^2=O(\lambda_C^2/L^2)$.
\item \textbf{Source's conformal factor in the universal metric (inconsistent).}
For $g=(1+\varepsilon)g_{\rm cl}$, the linearized Einstein tensor
$G_{00}=-\nabla^2\varepsilon$ carries no factor of $G$; for a Gaussian amplitude it
equals $6\lambda_C^2/L^4$. This contradicts F4 at order $\hbar^2G^0$, and would
require replacing F4.
\end{enumerate}
\end{theorem}

\noindent The result is a two-level master metric derived from the founding
principles:
\begin{equation}\label{eq:twolevel}
\boxed{g_{\mu\nu}=T^a{}_\mu T^b{}_\nu\Bigl[\eta_{ab}+\sum_ss_s\sigma^{(s)}_a\sigma^{(s)}_b+q_{ab}\Bigr],
\qquad \tilde g^{(\psi)}_{\mu\nu}=\Omega_\psi^2\,g_{\mu\nu},}
\end{equation}
with sector signs $s_s=-\varepsilon_s$ in the notation of Eq.~\eqref{eq:master}, and
where $q$ is fixed by F4 from the quantum stress of the source (vacuum
polarization, Eq.~\eqref{eq:ABquantum}, and the coherent-wave stress $T^Q$), and
$\Omega_\psi$ by Eq.~\eqref{eq:qcf}. This gives the founding master metric's quantum
slot, historically $-\lambda\mathcal Q$, a derived occupant. As $\hbar\to0$, $q$'s
quantum part and $\Omega_\psi-1$ vanish, recovering the classical restriction and
every classical result of Chapters~2--4. The structure reproduces quantum
mechanics in curved spacetime and semiclassical gravity, and it does not quantize
the metric: the limitations of semiclassical gravity, for example for macroscopic
superpositions of a source, apply unchanged.

\begin{proposition}[No natural-invariant dynamics beyond $g$ under F1--F5 (D)]\label{prop:nobeyond}
The microscopic Dirac action uses a vierbein of $g$ and is locally Lorentz invariant (Chapter~2,
Clifford compatibility), so it depends only on $g$. The Maxwell, scalar-contact and
Einstein--Hilbert terms and the quantum matter effective action $\Gamma[g]$ are
functionals of $g$. Theorem~\ref{thm:quantumno-go} therefore applies: dynamics
acting on more than $g$ requires an additional founding principle.
\end{proposition}

\subsection{The per-wave law for spin-$\frac12$ matter (QGD--112, QGD--113)}\label{sec:spinhalf}
For spin-$\frac12$ matter no metric conformal to $g$ makes the flow lines geodesic
(QGD--111). The exact replacement follows from the macroscopic spinor equation
alone. Define $\sigma=\bar\psi\psi\neq0$, the convective momentum
$P_\mu=-\operatorname{Im}(\bar\psi\nabla_\mu\psi)/\sigma$, the spin remainder
$\Lambda=\nabla\bar\psi\cdot\nabla\psi-|\bar\psi\nabla\psi|^2/\sigma$, the spin tensor
$S^{ab}=\bar\psi\Sigma^{ab}\psi$ with $\Sigma^{ab}=\frac i2[\gamma^a,\gamma^b]$, and the spin
vorticity $F=\mathrm dP$.
\begin{theorem}[Spin-$\tfrac12$ flow law in a curved universal metric (D)]\label{thm:spinhalf}
In Chapter~1's conventions, $(\gamma^\mu\nabla_\mu)^2=-\square+\frac14R$ and
$[\nabla_\mu,\nabla_\nu]=-\frac14R_{ab\mu\nu}\gamma^a\gamma^b$ on spinors. Consequently, exactly,
\begin{equation}\label{eq:spinhalfflow}
P\cdot P+M^2+\tfrac14R=\frac{\square\sqrt\sigma}{\sqrt\sigma}-\frac\Lambda\sigma,\qquad
F_{\mu\nu}=-\frac{\operatorname{Im}(\nabla_{[\mu}\bar\psi\nabla_{\nu]}\psi)}{\sigma}
-\frac{\partial_{[\mu}\sigma\,P_{\nu]}}{\sigma}-\frac{R_{ab\mu\nu}S^{ab}}{4\sigma},
\end{equation}
where $X_{[\mu\nu]}:=X_{\mu\nu}-X_{\nu\mu}$; and with $\tilde g=\Omega_{1/2}^2g$,
$\Omega_{1/2}^2=1+\frac{R}{4M^2}-\frac1{M^2}\bigl(\frac{\square\sqrt\sigma}{\sqrt\sigma}-\frac\Lambda\sigma\bigr)$
and $\tilde u_\mu=P_\mu/M$,
\begin{equation}\label{eq:spinhalflorentz}
\boxed{M\,\tilde u^\nu\tilde\nabla_\nu\tilde u_\mu=\tilde u^\nu F_{\nu\mu}.}
\end{equation}
\end{theorem}
\noindent The convective flow moves in the conformally rescaled metric like a
charge in the spin vorticity. $F$ contains a curvature--spin term with the
structure of the Mathisson--Papapetrou coupling, derived here rather than
imported. For spin-0 matter $\Lambda=0$ and $F=0$, and the law reduces to
Proposition~\ref{prop:qcf}.

\noindent\textbf{Verification.} Flat-region identities verified to $10^{-14}$ on exact
superpositions of plane-wave solutions (QGD--112). The curvature coefficients $+\frac14$
and $-\frac14$, the curved off-shell identity and the curved spin vorticity were verified
on an FLRW metric with an arbitrary spinor field, residuals $\le6\times10^{-16}$, with
the spinor connection constructed from $\nabla\gamma=0$ (QGD--113). In generic states the
spin-vorticity force is $40$--$70\%$ of the gradient force.

\noindent\textbf{Slow-motion limit (QGD--113).} With $\hat s$ the spin direction,
$\Lambda/\sigma\to\frac14\partial\hat s\cdot\partial\hat s$ (from the exact Fubini--Study identity
for unit two-spinors) and $F_{ij}\to-\frac12\hat s\cdot(\partial_i\hat s\times\partial_j\hat s)$, a
Mermin--Ho-type relation, both with $O(v^2)$ convergence verified. The kinetic
energy per particle is
$p^2/2M-(\hbar^2/2M)\nabla^2\sqrt n/\sqrt n+(\hbar^2/8M)|\nabla\hat s|^2$. For constant spin
the law reduces exactly to Chapter~1's irrotational result
$\dot{\bm p}=-\nabla(U+Q)$ and to the fixed-spin quantum stress of
Section~\ref{sec:quantumsource}. When spin varies, a spin quantum potential and a
spin-vorticity force appear that those fixed-spin treatments omit.

\noindent\textbf{Conditions.} $\sigma\neq0$ and $P$ timelike. As for spin-0 matter, the
law describes the convective momentum field; reading its integral curves as
trajectories is an interpretation.

\subsection{Where the natural invariant comes from}
\begin{proposition}[Closure identity (D)]\label{prop:closure}
On the static one-null-sector branch $g=\bar g+H\,\mathrm dv^2$ with $H\to0$ at
infinity, wherever $0<H<1$,
\begin{equation}\label{eq:closure}
\sigma_{\rm amp}=\sqrt H=\beta_{\rm ff}=\frac{2\zeta_{\rm ff}}{1+\zeta_{\rm ff}^2},\qquad
\zeta_{\rm ff}=\frac{\sqrt H}{1+\sqrt{1-H}},
\end{equation}
where $\beta_{\rm ff}$ is the speed, relative to static observers, of a test
particle released from rest at infinity, and $\zeta_{\rm ff}$ is its Chapter~1
on-shell spinor ratio. Covariantly, $\sigma_{\rm amp}^2=1-(u_{\rm ff}\cdot u_{\rm stat})^{-2}$.
The sector direction $k=\mathrm d(t+r)$ is a phase gradient solving the massless
limit of Chapter~1's Hamilton--Jacobi equation.
\end{proposition}
\noindent This extends Chapter~2's Schwarzschild diagnostic to the whole static
one-null-sector class. It fails for Kerr: the free-fall speed differs from the
Kerr--Schild amplitude, and the oblate $k$ is not closed. It also fails where
$FH\neq1$, for example in Eq.~\eqref{eq:ABquantum}. It establishes
self-consistency, not origin. The natural invariant's algebraic \emph{form} is
derived from the macroscopic spinor equation, while its exterior \emph{values}
for a given source are selected by F4.

\subsection{Lorentz covariance and gauge dependence}
The main equations were tested directly (QGD--110):
\begin{center}\small
\begin{tabular}{@{}p{0.45\linewidth}p{0.47\linewidth}@{}}\toprule
Equation or object & Result\\\midrule
Master metric, Eq.~\eqref{eq:master} & Invariant under frame Lorentz transformations ($T\mapsto\Lambda T$, $\sigma_a\mapsto\sigma_c(\Lambda^{-1})^c{}_a$); tensorial under coordinate changes\\
Historical $\mathcal Q$ with $\bar\nabla$ & Tensorial: Cartesian and spherical computations agree exactly\\
Macroscopic spinor equation (Ch.~1 conventions) & Covariant: a spinor representation $S(\Lambda)$ exists and $S u(p)$ solves the equation at $\Lambda p$\\
Spinor ratio $\mathsf R(p)$ & Transforms exactly as $\mathsf R(p)\mapsto\mathsf R(\Lambda p)$\\
$\zeta=p/(E+Mc^2)$ & Poincar\'e-ball coordinate of the mass shell; boosts act by M\"obius addition (exact). $\beta=2\zeta/(1+\zeta^2)$ is the Beltrami--Klein coordinate\\
Spinor component moduli, $\psi^\dagger\psi$ & \textbf{Not} invariant; $\bar\psi\psi$ and $j^\mu j_\mu$ are\\
$\Omega_\psi$, $T^Q$ (scalar wave) & Covariant\\
Sector representation $\sigma$ & \textbf{Gauge dependent}: the same Schwarzschild geometry (Kretschmann $12r_s^2/r^6$ in both) has $h=g-\bar g$ of rank 1 in the advanced identification and rank 2 in the static one\\\bottomrule
\end{tabular}
\end{center}
\noindent Three consequences follow:
\begin{enumerate}
\item \textbf{Spin-$\frac12$ version of the quantum terms (QGD--111--113).} The scalar-wave
Eqs.~\eqref{eq:qcf} and~\eqref{eq:TQ} use component amplitudes, which are not
invariant for spinors. The covariant spin-$\frac12$ amplitude is $(-j\cdot j)^{1/4}$
(Fierz identity, verified). The per-wave level for spin-$\frac12$ matter is
Theorem~\ref{thm:spinhalf}. The universal level holds unchanged: in the energy frame,
the stress tensor of a coherent spin-$\frac12$ field is $\rho\,u_Lu_L+\Pi$ with
$\Pi=O((\lambda_C/L)^2)\rho$ (numerical scaling, slope 1.98), and the probability
current tilts relative to the energy flow at $O(\lambda_C/L)$ (slope 1.01) through the
spin current, which is not an extra stress.
\item \textbf{Geometric meaning of the spinor ratio.} The Chapter~1 ratio $\zeta$
is not a Lorentz vector. It is a hyperbolic-geometry coordinate (Poincar\'e
ball), and the closure identity~\eqref{eq:closure} is the ball-to-Klein map.
\item \textbf{The natural invariant is not a diffeomorphism invariant.}
Because the master metric contains a fixed reference $\bar g$, the sector
representation, including the number of sectors, their rank and their signs,
depends on how spacetime is identified with the reference. Only $g$, and
quantities built from it, are gauge invariant. Every QGD prediction under F4 is
computed from $g$, so no prediction is affected, but statements about $\sigma$
must state the identification used.
\end{enumerate}
% =================== END QGD-110 INSERTION ===================

\section{What the equation makes possible---and what it does not}
\label{sec:physics}
Equation~\eqref{eq:split} has an important explanatory role.
It expresses effective gravitational equilibrium as a balance of the
metric-mediated gravitational action against physical matter stress,
projected into the natural-invariant variables. The source is the
\emph{complete, conserved} stress of the physical system; a local
Dirac-component ratio or a single wave amplitude is not generally an
adequate universal source. The Dirac state may contribute through
basis-independent currents, energies and stress correlations.

But the action-derived equation also establishes a \emph{limit}:
with no direct natural-invariant interactions and no new dynamical
operator, the effective field equations contain no distinct
propagating $\sigma$ mode beyond the effective metric degrees of
freedom. A new observable could follow from (i) a well-defined direct
coupling between quantum matter and natural-invariant data, (ii) a
physically derived independent natural-invariant kinetic/constraint
sector, or (iii) a genuinely distinct quantum effective action.
None follows inevitably from the signed metric identity; each would
require a separately specified covariant action, conservation laws,
correct field content and measurable consequences.

% === QGD-101 insertion, corrected QGD-105 ===
\paragraph{Note on route (i) (QGD--100, corrected QGD--105).}
Research note QGD--100 examined one \emph{hypothetical} realization of route~(i):
removing the reference term, reading four $\eta$-signed sectors as a tetrad, and
assuming a microscopic theory with only global frame symmetry. Neither assumption
belongs to the formulation of Chapters~1--4. Under them the two-derivative
dynamics is New General Relativity \cite{HS}, whose parameter directions other than the
axial one give $\gamma=-2c/(a+b-4c)\ne1$ at linear order, and whose axial
direction is obstructed at nonlinear order according to \cite{BJD}.
This rules out that particular hypothetical route; it says nothing about
routes (i)--(iii) in general.
% === end QGD-101/105 insertion ===

A particularly sharp discriminator is the comparison of different
physical quantum source states with matched mean stress and classical
boundary data. A measured mean-field difference beyond standard
GR plus quantum-field corrections would indicate a new physical
coupling. Merely assigning different signed representatives to the
same metric cannot change a metric-only observable.

\section{Relationship to the archived historical Chapter 4}
\label{sec:historical}
The historical research program wrote schematically
\begin{equation}\label{eq:historical}
\square_g\sigma_\mu=Q_\mu+G_\mu+T_\mu
+\kappa\ell_Q^2\square_g^2\sigma_\mu+\cdots.
\end{equation}
This is \emph{not} equivalent by definition to
Eq.~\eqref{eq:split}. The historical $Q_\mu,G_\mu,T_\mu$
expressions have unresolved algebraic/provenance and normalization
issues documented in the QGD--052 audit. In particular, no term
$\square_g\sigma_\mu$ or $\square_g^2\sigma_\mu$
can be deduced simply from varying the derivative-free composite
Einstein--Hilbert action. Adding a standalone $\sigma$ kinetic term
would be a change of the effective theory and requires physical
justification, not a cosmetic rewriting of the balance equation.
The archived manuscript therefore remains an independent historical
source, while the present chapter begins with the precise equation
supported by the accepted action and metric map.

\section{Conclusion and the next legitimate question}
The action-based content of this updated Chapter~4 has one central field-equation result:
\begin{equation}\label{eq:conclusion}
\boxed{\cS_{(s)}^\rho
=2s_s\cE_g^{\rho\nu}\sigma_\nu^{(s)}=0,\qquad
\cS_{(s)}^\rho=\cG_{(s)}^\rho+\cT_{(s)}^\rho+\cQ_{(s)}^\rho.}
\end{equation}
For the metric-only effective action, $\cQ=0$ and, on the
full-rank signed branch, this balance reproduces the established
Einsteinian field equations. This is a rigorous mathematical bridge
from QGD's natural-invariant variables to effective gravity, not a
claim that the underlying physical origin of gravity has been solved.
The question that can make QGD more than a geometric reformulation
is therefore precise: \emph{Does the quantum matter-wave structure
lead to independent, covariantly defined natural-invariant dynamics
or source dependence, and does any such effect produce an observable
not already predicted by GR and quantum field theory?}
\par\medskip
The Newton-calibrated Chapter~3 potential and the controlled $\hbar^2$
quantum-source stress are compatible with this balance, but neither
supplies an independent natural-invariant propagation operator.
The metric-only quantum-equivalence theorem in
Section~\ref{sec:independent} makes this limitation exact under
its stated rank and action assumptions; the historical attractive
$9/2$ force would additionally require a nonzero exterior source
residual or a changed gravitational response. The
historical $9/2$ correction remains unverified. Future extensions must
derive their covariant operators, coefficient matching and physical
observables before assigning a genuinely new wave equation or degree
of freedom to $\sigma_\mu$.


% === QGD-101 insertion (moved from the abstract) ===
\section*{Revision history}\label{sec:revision-history}
\addcontentsline{toc}{section}{Revision history}
\textbf{QGD--096} established the equivalence/no-independent-correction
theorem for the regular metric-only theory and calculated the nonzero
vacuum-Poisson residual that the historical attractive $9/2$ force would
require. \textbf{QGD--097} added a separately sourced semiclassical
benchmark for a massless four-component Dirac vacuum field: a verified
far-field $-2G_N^2M\hbar/(15\pi c^3r^3)$ correction to the static
Killing-energy potential, taken from published quantum-matter backreaction,
not from the historical complex-phase Taylor expansion or a new QGD
operator. \textbf{QGD--098} independently derived and checked the
angular-channel coefficients of the Schwarzschild radial Dirac reduction and
recorded why its Liouville-normal effective potential cannot be promoted to
a physical force by differentiation. \textbf{QGD--101} integrated research
notes QGD--099 and QGD--100. \textbf{QGD--105} audited those insertions:
Remark~\ref{rem:Q-absorbed} was restricted to Chapter~2's constructed branch;
the route-(i) paragraph was restated as a test of one hypothetical realization;
and a closing paragraph that reframed QGD's identity was removed, because it
did not follow from the chapter's stated assumptions. No earlier equation or
label was removed. \textbf{QGD--110} added Section~\ref{sec:quantumstructure}: the
split identity, the quantum term in static variables, the test of the historical
derivative trial, the derived quantum conformal factor, the two-level master metric,
the closure identity, and Lorentz-covariance and gauge tests. \textbf{QGD--111} replaced
the open spin-$\frac12$ item with covariant results; \textbf{QGD--112} added the exact
spin-$\frac12$ flow law and restricted the name ``Dirac'' to the microscopic field.
\textbf{QGD--113} added the curved-metric form of that law (Theorem~\ref{thm:spinhalf}),
its slow-motion limits, and the cross-reference in Section~\ref{sec:quantumsource}.
\textbf{QGD--115} added Section~\ref{sec:lovelock}: the Einstein--Hilbert form derived by
uniqueness under Principle~L, with the coupling from the Newtonian calibration.
% === end QGD-101 insertion ===

\section*{Source and research provenance}
\begin{thebibliography}{9}
\bibitem{ch1} QGD, \emph{Chapter 1: Foundations}, accepted Dirac/matter-wave construction (existing research manuscript).
\bibitem{ch2} QGD, \emph{Chapter 2: Field Equations}, complete signed master metric, variational right-inverse and effective action (QGD--075).
\bibitem{ch3} QGD, \emph{Chapter 3: Gravitational Normalization and Quantum-Force Verification} (QGD--087 surgical revision; QGD--101 notation update).
\bibitem{audit} QGD--052, \emph{Focused Historical QGT/QGD Field-Equation Provenance Audit}; and QGD--094, \emph{General Natural-Invariant Field Equations and Historical Source Split}, working research checks.
\bibitem{referee} QGD--091--092, \emph{Independent Schwarzschild Dirac radial reduction and referee verification}, conditional test-field coefficient audit (cross-checked in QGD--098).
\bibitem{quantum} QGD--090--094, \emph{Quantum Force, Dirac Radial Dynamics, and Quantum-Source Verification}, conditional source and effective-metric studies. Historical $\hbar^2$ operators remain unmatched.
\bibitem{Carlson2010} E. D. Carlson, P. R. Anderson, A. Fabbri, S. Fagnocchi, W. H. Hirsch, and S. A. Klyap, ``Semiclassical gravity in the far field limit of stars, black holes, and wormholes,'' \emph{Phys. Rev. D} \textbf{82}, 124070 (2010), DOI: 10.1103/PhysRevD.82.124070; see Eqs. (7.2), (8.3), (8.6), and (8.7).
\bibitem{donoghue} J. F. Donoghue, ``General Relativity as an Effective Field Theory: The Leading Quantum Corrections, \emph{Phys. Rev. D} \textbf{50}, 3874 (1994), DOI: 10.1103/PhysRevD.50.3874.
\bibitem{equivalence} S. Kamefuchi, L. O\textquoteright Raifeartaigh, and A. Salam, ``Change of variables and equivalence theorems in quantum field theories, \emph{Nuclear Physics} \textbf{28}, 529--549 (1961), DOI: 10.1016/0029-5582(61)91075-6.
\bibitem{qgd100} QGD--100, \emph{Testing the Global-Lorentz Branch}, research note (2026).
\bibitem{HS} K.~Hayashi and T.~Shirafuji, ``New general relativity,'' \emph{Phys. Rev. D} \textbf{19}, 3524 (1979).
\bibitem{Lovelock1971} D.~Lovelock, ``The Einstein tensor and its generalizations,'' \emph{J. Math. Phys.} \textbf{12}, 498 (1971).
\bibitem{Lovelock1972} D.~Lovelock, ``The four-dimensionality of space and the Einstein tensor,'' \emph{J. Math. Phys.} \textbf{13}, 874 (1972).
\bibitem{Jacobson} T.~Jacobson, ``Thermodynamics of spacetime: the Einstein equation of state,'' \emph{Phys. Rev. Lett.} \textbf{75}, 1260 (1995).
\bibitem{Sakharov} A.~D.~Sakharov, \emph{Sov. Phys. Dokl.} \textbf{12}, 1040 (1968).
\bibitem{Visser} M.~Visser, ``Sakharov's induced gravity: a modern perspective,'' \emph{Mod. Phys. Lett. A} \textbf{17}, 977 (2002).
\bibitem{BJD} J.~Beltr\'an Jim\'enez and K.~F.~Dialektopoulos, ``Non-linear obstructions for consistent new general relativity,'' \emph{JCAP} \textbf{2020}(01), 018, arXiv:1907.10038.
\end{thebibliography}
\end{document}
